\chapter{CHAPTER 1: INTRODUCTION}\label{chapter-1-introduction}

\section{1.1 Background and Motivation}\label{background-and-motivation}

Air pollution is increasingly recognized as a critical global health
crisis. According to the World Health Organization (WHO), prolonged
exposure to fine particulate matter (PM2.5) is linked to severe
cardiovascular and respiratory diseases, accounting for millions of
premature deaths globally (WHO, 2021). In South Asia, rapid urbanization
and industrial growth have exacerbated air quality deterioration. Sri
Lanka, while historically enjoying better air quality than its
continental neighbors, has experienced significant pollution spikes in
recent years. These spikes are primarily driven by vehicular emissions,
industrial activities, and transboundary air pollution carried by
seasonal monsoonal winds (Rathnayake et al., 2023).

In Sri Lanka, the two most economically and culturally significant urban
centers are Colombo and Kandy. These cities possess distinct
topographies that create unique microclimates and dispersion mechanisms.
Colombo, the commercial capital located on the western coastal plain,
typically benefits from the ``sea breeze effect'' which aids in
pollution dispersion, yet suffers from high chronic vehicular emissions.
Conversely, Kandy is a heritage city situated within a central highland
valley. This geography creates a ``valley trapping'' effect, where
surrounding mountains prevent the dispersion of locally generated
emissions, particularly during frequent thermal inversions (Ambanwala,
2025).

Despite these growing public health risks, Sri Lanka currently relies on
sparse, retrospective data networks---such as the low-cost PurpleAir
sensor network---which report current conditions but lack predictive
capabilities (Dhammapala et al., 2022). To enact proactive health
management and targeted interventions (e.g., closing schools or issuing
traffic restrictions), citizens and municipal authorities require highly
accurate, localized forecasts rather than reactive measurements.

\section{1.2 Problem Statement}\label{problem-statement}

The primary problem addressed in this research is the absence of a
reliable, localized, and multi-horizon PM2.5 forecasting system in Sri
Lanka that structurally accounts for the country's distinct urban
topographies. Existing monitoring solutions merely report historical or
present conditions, leaving a critical gap in proactive, predictive
awareness. Furthermore, standard high-accuracy machine learning
forecasting models operate as ``black boxes,'' providing numerical
predictions without elucidating the underlying drivers of the pollution
(e.g., whether an impending spike is driven by historical accumulation,
humidity, or cyclical time patterns). This lack of model transparency
severely hinders the ability of environmental authorities to trust the
predictions and implement effective, targeted public health
interventions.

\section{1.3 Research Gap}\label{research-gap}

Current literature concerning air quality forecasting in tropical,
developing nations reveals significant limitations, specifically in the
Sri Lankan context. This research explicitly addresses three identified
gaps:

\begin{enumerate}
\def\labelenumi{\arabic{enumi}.}
\tightlist
\item
  \textbf{The Multi-Horizon Predictive Gap:} What is missing is a
  comprehensive benchmarking of diverse algorithmic families across
  extended forecasting windows. While some local studies (Kurukulasuriya
  et al., 2025; Fernando et al., 2022) have explored short-term air
  quality trends using isolated models, there is no systematic
  evaluation comparing statistical (SARIMAX), machine learning
  (XGBoost), and deep learning (LSTM) approaches across 1, 6, 12, 24,
  and 48-hour horizons.
\item
  \textbf{The Topographical Comparative Gap:} Existing predictive models
  in the region are predominantly trained on homogeneous, coastal
  terrains (Colombo) and lack spatial awareness. What hasn't been done
  is a direct, empirical comparison of how forecasting accuracy and
  optimal feature sets degrade or shift when applied to a valley basin
  (Kandy) versus a coastal metropolis (Colombo).
\item
  \textbf{The Interpretability Gap:} High-accuracy predictive models in
  the local literature universally suffer from the ``black box''
  trade-off (Attanayake et al., 2025). There is a complete absence of
  Explainable AI (XAI) frameworks---specifically SHAP (SHapley Additive
  exPlanations)---applied to Sri Lankan air quality forecasting to
  quantify and interpret exactly \emph{why} a model predicts a severe
  PM2.5 event.
\end{enumerate}

\section{1.4 Research Aim}\label{research-aim}

To \textbf{design} a multi-horizon, feature-engineered machine learning
forecasting framework integrated with Explainable AI (SHAP),
\textbf{develop} an automated inference pipeline culminating in a
consumer-facing mobile application, and \textbf{evaluate} its predictive
performance and interpretability across the distinct coastal and valley
topographies of Colombo and Kandy.

\section{1.5 Research Objectives}\label{research-objectives}

To achieve the stated aim, the following Specific, Measurable,
Achievable, Relevant, and Time-bound (SMART) objectives have been
formulated:

\begin{enumerate}
\def\labelenumi{\arabic{enumi}.}
\tightlist
\item
  To \textbf{collect and synchronize} a multi-modal dataset comprising
  hourly PurpleAir PM2.5 ground telemetry and Open-Meteo meteorological
  data for Colombo and Kandy, spanning from January 2022 to December
  2024.
\item
  To \textbf{engineer} a 48-dimensional feature vector---incorporating
  complex time-lagged variables, rolling statistical means, and cyclical
  trigonometric time encodings---to capture non-linear atmospheric
  patterns.
\item
  To \textbf{design and train} a statistical baseline (SARIMAX), an
  ensemble machine learning model (XGBoost), and a sequential deep
  learning architecture (LSTM Encoder-Decoder) for PM2.5 forecasting.
\item
  To \textbf{evaluate} model accuracy across 1, 6, 12, 24, and 48-hour
  horizons on unseen 2025 test data using Root Mean Squared Error (RMSE)
  and Mean Absolute Error (MAE), explicitly comparing performance
  between coastal and valley topographies.
\item
  To \textbf{implement} Explainable AI (SHAP) to quantify, visualize,
  and interpret the top three meteorological and temporal features
  driving predictions for both cities.
\item
  To \textbf{develop} an automated backend Python inference pipeline
  that fetches real-time data, executes the optimal forecasting model,
  and writes the resulting predictions to a Firebase cloud database.
\item
  To \textbf{design and deploy} a React Native mobile application
  (SentinelAQ) that consumes the Firebase payload, providing end-users
  with real-time AQI dials, 48-hour charts, SHAP insights, and
  threshold-based push notifications.
\end{enumerate}

\section{1.6 Research Challenges}\label{research-challenges}

Addressing the identified problem domain presents several notable
challenges: * \textbf{Data Quality and the Pivot from Satellite Fusion:}
The initial architecture proposed the spatio-temporal fusion of daily
Sentinel-5P satellite chemical vectors (NO2) with hourly ground data.
However, a major domain-specific challenge was encountered: severe cloud
occlusion over Kandy's valley and Sri Lanka's tropical monsoons resulted
in massive satellite data gaps. Furthermore, the frequency mismatch
(daily satellite passes vs.~hourly ground fluctuations) degraded
short-term forecasting accuracy. This necessitated a strategic
methodological pivot to rely on robust, high-frequency meteorological
proxy data (Open-Meteo) rather than satellite imagery. * \textbf{The
``Black Box'' Trade-off:} Deep learning and ensemble models inherently
lack transparency. Integrating SHAP is computationally expensive and
complex to implement for multi-horizon regression tasks, but it is an
absolute necessity to secure end-user trust and policy support. *
\textbf{Multi-Horizon Error Propagation:} Forecasting accuracy naturally
and exponentially degrades over extended horizons (e.g., 48 hours). The
technical challenge lies in engineering sufficiently robust
autoregressive features (like 24-hour rolling means) that allow the
models to maintain acceptable error margins two days into the future.

\section{1.7 Contributions}\label{contributions}

This research provides the following novel contributions to knowledge,
methodology, and practice: * \textbf{Methodological Innovation:} It
provides empirical validation that feature-engineered, tree-based
ensemble learning (XGBoost) significantly outperforms complex sequential
deep learning (LSTM) for multi-horizon environmental forecasting in
tropical, developing-nation microclimates. * \textbf{Technical
Contribution:} This study marks the first application of SHAP
(Explainable AI) to multi-horizon PM2.5 forecasting in Sri Lanka. It
pioneers a transparent approach that explicitly quantifies how
topographical differences alter the driving factors of pollution (e.g.,
local accumulation in Kandy versus cyclic weather dependencies in
Colombo). * \textbf{Practical Impact:} It delivers ``SentinelAQ,'' a
fully functional, end-to-end prototype---from automated data ingestion
to an intuitive mobile application. This demonstrates a scalable,
cost-effective framework for deploying predictive environmental AI
directly to citizens, circumventing the need for expensive government
sensing infrastructure.

\section{1.8 Chapter Summary}\label{chapter-summary}

This chapter established the critical need for a predictive,
transparent, and multi-horizon PM2.5 forecasting system tailored to Sri
Lanka's unique urban topographies. By identifying the limitations in
current retrospective monitoring and the ``black box'' nature of
existing predictive models, explicit research gaps were defined. The
subsequent aim and SMART objectives laid out a structured framework to
design, develop, and evaluate the SentinelAQ system, while acknowledging
the inherent meteorological and technical challenges. The following
chapter provides a comprehensive review of related literature,
critically evaluating existing approaches to air quality forecasting and
establishing why the methodologies selected for this study are
necessary.
